% =====================================================================
%  IMAGEN — poner una imagen sola, centrada.
%
%  \imagen{nube.png}                 <- ancho automatico (\linewidth)
%  \imagen[0.5\linewidth]{nube.png}  <- ancho a medida
%
%  El nombre puede ir con espacios: \imagen{ nube.png } funciona.
%
%  DONDE BUSCA LA IMAGEN (en este orden):
%
%    1. La ruta tal cual, desde la carpeta donde compilas.
%       \imagen{Tema1.1/imagenes/nube.png}
%
%    2. La carpeta "imagenes_explicativas" del proyecto (y sus
%       subcarpetas).  Se busca por el NOMBRE del fichero, asi que da
%       igual en que subcarpeta este:
%       \imagen{nube.png}  ->  imagenes_explicativas/tus-apuntes/nube.png
%
%    3. Si no aparece en ningun sitio, se busca por el nombre en todo
%       el proyecto (carpeta actual y las de arriba).
%
%  Pon siempre la extension del fichero (nube.png, nube.jpg...).
%  Las rutas se cuentan desde la RAIZ del documento (la carpeta donde
%  compilas). Necesita -shell-escape (ya activado en VS Code).
% =====================================================================

% macro vacia para comparar con \ifx (\@empty no sirve aqui)
\def\imagenvacio{}

% ---------------------------------------------------------------------
%  Estas dos van en modo normal (fuera de \ExplSyntaxOn) porque
%  necesitan espacios de verdad: el comando del shell y el mensaje
%  de error.
% ---------------------------------------------------------------------
%  #1 = carpeta donde se deja la copia
%  #2 = nombre del fichero que se busca
%  #3 = ruta del fichero copiado
%  #4 = ruta del script que hace la busqueda
\newcommand{\imagenbuscar}[4]{%
  \immediate\write18{mkdir -p "#1" && sh "#4" "#2" "#3"}}

\newcommand{\imagenerror}[1]{%
  \PackageError{imagen}{No se encuentra la imagen "#1"}{%
    Ponla en la carpeta "imagenes_explicativas" del proyecto (por ejemplo
    imagenes_explicativas/tus-apuntes/) o escribe la ruta completa desde
    la raiz del documento. No olvides la extension del fichero (por
    ejemplo .png).}}

% ---------------------------------------------------------------------
%  \buscarimagen{fichero}
%
%  Busca la imagen y deja la ruta buena en \imagenruta.
%  La usan tambien \filaimagen y \filaimagender.
%
%  OJO: \IfFileExists solo entiende el nombre si es UN solo macro (o
%  texto plano).  Por eso las rutas se montan antes en \imagentmpdir
%  y \imagentmpfile, que son solo caracteres.
%
%  OJO 2: la busqueda la hace el script buscarimagen.sh, porque en
%  TeX el ";" es un caracter de comentario y un \write18 con ";" se
%  cortaba a la mitad.
% ---------------------------------------------------------------------
\ExplSyntaxOn

\tl_new:N \l__imagen_nombre_tl

\cs_new_protected:Npn \buscarimagen #1
  {
    % --- el nombre del fichero, sin espacios en los extremos ---
    \tl_set:Nn \l__imagen_nombre_tl { #1 }
    \tl_trim_spaces:N \l__imagen_nombre_tl
    % --- rutas "escapadas" (solo caracteres) ---
    \edef\imagenficheroesc { \tl_to_str:N \l__imagen_nombre_tl }
    \edef\imagenjobesc { \expandafter\detokenize\expandafter{ \jobname } }
    \edef\imagentmpdir { /tmp/apuntes-imagen-\imagenjobesc }
    \edef\imagentmpfile { \imagentmpdir/\imagenficheroesc }
    % --- 1) la ruta tal cual, sin buscar ---
    \IfFileExists { \imagenficheroesc }
      { \edef\imagenruta { \imagenficheroesc } }
      {
        % --- 2) "imagenes_explicativas" del proyecto (y subcarpetas) ---
        % --- 3) y si no, por el nombre en todo el proyecto ---
        \imagenbuscar
          { \imagentmpdir }
          { \imagenficheroesc }
          { \imagentmpfile }
          { \bibliotecasDir/estilo/filas/buscarimagen.sh }
        \IfFileExists { \imagentmpfile }
          { \edef\imagenruta { \imagentmpfile } }
          {
            % --- 4) no esta en ninguna parte ---
            \edef\imagenruta { \imagenficheroesc }
            \imagenerror { \imagenficheroesc }
          }
      }
  }

\ExplSyntaxOff

% ---------------------------------------------------------------------
%  \imagen[ancho]{fichero}
% ---------------------------------------------------------------------
\newcommand{\imagen}[2][]{%
  \begingroup
    % --- ancho (y alto maximo, para que la imagen no se coma la pagina) ---
    \def\imagenancho{#1}%
    \ifx\imagenancho\imagenvacio\def\imagenancho{\linewidth}\fi
    \def\imagenalto{0.68\textheight}%
    % --- buscar la imagen ---
    \buscarimagen{#2}%
    % --- meterla ---
    \IfFileExists{\imagenruta}{%
      \begin{center}%
        % graphicx no expande el nombre del fichero: hay que pasarlo
        % "escapado" (detokenize) para que los _ y otros signos valgan.
        \edef\imagenrutadetok{\expandafter\detokenize\expandafter{\imagenruta}}%
        \edef\imagenllamada{\noexpand\includegraphics[width=\imagenancho,height=\imagenalto,keepaspectratio]{\imagenrutadetok}}%
        \imagenllamada
      \end{center}%
    }{}%
  \endgroup
}
